\scanpage{35}
Explicitons l'axiome

D) Toute \deleted{\unsure{ens}} \inserted{partie} minorée \inserted{$\ne\varnothing$} \deleted{\unsure{d'éléments}} \inserted{de $\mathcal X_{\unsure{**|xx}}=\mathcal X_0\cup\mathcal X_1\cup\mathcal X_2$} admet un inf,
($\Longleftrightarrow$ toute partie majorée \inserted{$\ne\varnothing$} de $\mathcal X_{**}$ admet un sup).

\deleted{Comme il résulte de c) que les él. minimaux de $\mathcal X_{**}$ sont
les él. de $S$, les parties minorées de $\mathcal X_{**}$ sont celles
contenues dans un $\mathcal X_{\geq s}$. Or comme ce dernier
c'est un polygône \inserted{(avec plus petit él.)}, et dans un \inserted{tel} polygône \unsure{satisfaire|satisfaire à|satisfait} la
condition \inserted{d'existence} de borne inférieure \unsure{sss|ssi|ses} \deleted{\unsure{il a au}} son nb de
côtés est $\geq3$ (dans le cas excluant celui $n=2$ ici\ldots)
Donc la condition D se récrit :

d) $\forall s\in S$, il y a au moins trois arêtes \inserted{$a\in A$} incidentes à $s$.}

Il suffit d'exprimer que deux él. distincts \inserted{minorisés \unsure{égalts|dgts|distincts}} de $\mathcal X_{**}$ \deleted{\unsure{ont}} \inserted{majorant un él. de \unsure{$S$}} un inf. Si l'un des deux est $\in\mathcal X_0$ (i.e. est minimum)
c'est trivialement satisfait. Il reste à traiter $3$ cas

1°) \uline{Deux arêtes} : on trouve que 2 arêtes qui ont \unsure{ms|m\textsuperscript{es}}
sommets sont égales, i.e. $(S,A,\vec A)$ est un graphe strict.

2°) \uline{Deux faces} : on trouve \deleted{\unsure{deux faces qui|qu'elles} $(F,A,A^\uparrow)$} \inserted{que les deux faces distinctes}
ont en commun deux sommets, \deleted{\unsure{elles|elles-mêmes}} ceux-ci sont liés
\deleted{\unsure{par une arête qui leur soient}} et elles \deleted{\unsure{ont}} ont en commun l'arête qui les joint,
\inserted{et pas d'autre sommet}

3°) Une arête \inserted{$a$} et une face \inserted{$f$} \deleted{non} \inserted{une arête} incidente --- mais
ayant un sommet commun : \deleted{il faut} exprimer
que le deuxième sommet de $a$ n'est pas dans $f$,
ou encore \deleted{que} : si $f$ est incidente aux
deux sommets d'une arête, elle est incidente
à l'arête.

Regroupant, on trouve

\margin{toute arête est le sup de l'ens de ses 2 extrémités}
$d_1)$ Une arête est déterminée par \deleted{ses} \inserted{l'ens de ses} deux extrémités

$d_{\unsure{2}})$ Une face qui cont les 2 extr. d'une arête contient celle-ci

\margin{toute arête est le Inf de l'ens des deux faces incidentes}
$\unsure{d_1|d'_1})$ Une arête est déterminée par l'ens des deux
faces incidentes

$d'_2)$ Les deux faces incidentes à une arête n'ont d'autre
sommet incident commun que ses extrémités
% archive sheet number: 34
